% page 62 du fac-simile (image p-061 du scan)
% bloc 160.0 x 237.7 mm ; marge gauche 8.0 mm
\pgc{-12.700mm}{0.000mm}{
\l{\cel{6}54}
\l{}
\l{\sou{bakteriologio}. La parto di la biologio, qua havas kom temo}
\l{\cel{3}di studio la bakterii. - DEFIRS.}
\l{}
\l{\sou{bakteriologo}.\cel{2}Ciencisto pri bakteriologio. - DEFIRS.}
\l{}
\l{\sou{balo}.\cel{2}Asemblitaro grandanombra de personi qui dansas. - DEFIRS}
\l{}
\l{\sou{balado}. Mikra poemo ek versi inter-egala, qua konsistas ek}
\l{\cel{3}tri strofi simetra ed un kupleto uzata kom homajo-versi.}
\l{\cel{3}- DEFIRS.}
\l{}
\l{\sou{balanco}. Instrumento por ponderar korpo, equilibriganta}
\l{\cel{3}olca per maso de metalo uzata kom pez-unajo. - DEFIRS.}
\l{}
\l{\sou{balanciero}. Maso (generale roto-forma) qua ocilas cirkum}
\l{\cel{3}sua baricentro per resorto. - DEFIRS.}
\l{}
\l{\sou{balas-rubino}. Lapido ek aluminato di magnezio, varietato}
\l{\cel{3}rozea. - DEFIRS.}
\l{}
\l{\sou{balasto}. Grosa sabl-uni o fer-skorio, uzata : l. en la mar-}
\l{\cel{3}navigado, por igar la navi plu equilibroza, per abasar}
\l{\cel{3}la baricentro ; 2. sur la relvoyi, por kovrar la traversi}
\l{\cel{3}qui suportas la reli. - DEFRS.}
\l{}
\l{\sou{balayar}.\cel{2}(\sou{trans, per}). Pulsar de avan su, sur sulo - polv-}
\l{\cel{3}uni, rezidui, aquo, nivo-floki, e c., per fasko de mikra}
\l{\cel{3}stipi o per brosilo longa-krina, fixigita a mancho longa.}
\l{}
\l{}
\l{\sou{balbutar}.\cel{2}(\sou{trans.}) Artikular la vorti en maniero poke di-}
\l{\cel{3}cernebla. - FIS.}
\l{}
\l{\sou{balde}. Poka tempo pose. - D.}
\l{}
\l{\sou{baldakino}.- I. Doselo, subtenata da koloni, e garnisita ye}
\l{\cel{3}tapeti, qua aspektas quale krono super la altaro, en la}
\l{\cel{3}kirki. - II. Lito-doselo de ube pendas kurteni. - DEFIRS.}
\l{}
\l{\sou{baldrio}. Larja bendo de ledro o de stofo qua, suportante la}
\l{\cel{3}espado o la sabro, portesas da la persono diagonale, de}
\l{\cel{3}la shultro dextra a la hancho sinistra. - DEFI.}
\l{}
\l{\sou{baleno}. (\sou{zool.)}\cel{2}Granda mamifero, ek la klaso \textquotedbl{}cetacei\textquotedbl{}, di}
\l{\cel{3}qua la boko sen-denta, esas garnisita - ye la du lateri,}
\l{\cel{3}an la maxilo) ye granda lami dina, flexebla (barti).}
\l{\cel{3}- L.\cel{1}\sou{balaena}. - FISL.}
\l{}
\l{\sou{baleto}.\cel{2}Danso alegoria, exekutata sive en festo, sive en}
\l{\cel{3}teatro, kom intermezo di opero, di komedio. - DEFIRS.}
\l{}
\l{\sou{balifo}. Oficiro nobela o yurala qua, en Francia, olim judi-}
\l{\cel{3}ciis ye la nomo di la rejo o di sinioro. - EFIS.}
}
